%%  AlphaZero Gomoku – Comprehensive Study Report
%%  compile:  pdflatex report.tex   (or generate_report.py → report_final.tex)

\documentclass[11pt,a4paper]{article}

\usepackage[margin=2.5cm]{geometry}
\usepackage{amsmath,amssymb}
\usepackage{booktabs}
\usepackage{graphicx}
\usepackage{xcolor}
\usepackage{hyperref}
\usepackage{caption}
\usepackage{subcaption}
\usepackage{listings}
\usepackage{algorithm}
\usepackage{algpseudocode}
\usepackage{multirow}
\usepackage{array}
\usepackage[round]{natbib}

\hypersetup{
  colorlinks=true, linkcolor=blue!70!black,
  citecolor=green!50!black, urlcolor=blue!70!black,
}

\lstset{
  basicstyle=\ttfamily\small,
  frame=single,
  breaklines=true,
  numbers=left,
  numberstyle=\tiny\color{gray},
}

\title{%
  \textbf{AlphaZero-Style Gomoku: Architecture, Training Dynamics,\\
  and Strategic Intelligence across Five Configuration Studies}
}
\author{Mehul Chourasia\\
  \small UCL Department of Computer Science\\
  \small Reinforcement Learning Coursework 2025--26
}
\date{\today}

% ============================================================
\begin{document}
\maketitle

\begin{abstract}
We present a comprehensive empirical study of an AlphaZero-style reinforcement
learning agent applied to $9{\times}9$ Gomoku.  The agent combines Monte Carlo
Tree Search (MCTS) guided by a policy-value neural network with self-play data
generation, board-symmetry augmentation, and an optional heuristic-opponent
curriculum.  Five distinct configurations are trained for 250 iterations apiece
(\emph{in parallel, from scratch}), varying network depth (5–10 residual blocks),
channel width (64–128), MCTS simulation budget (200 fixed; $100 \to 200 \to 400$
progressive), learning-rate schedule (MultiStep vs.\ Cosine annealing), and
gradient clipping.  We report learning curves, ablation results, and head-to-head
tournament outcomes, and complement these with a ten-scenario strategic evaluation
battery that probes the agent's tactical capabilities (winning threats, blocking,
forks, and counter-fork play).  The best configuration (\textbf{%%BEST_CFG%%})
achieves a win rate of %%BEST_WR_HEURISTIC%% against the threat-based heuristic
opponent, while tournament analysis reveals that larger networks with progressive
simulation schedules dominate smaller, fixed-sim baselines.  Strategic evaluation
further confirms that all mature checkpoints handle immediate threats reliably,
while fork creation and counter-fork play only emerge after $\sim150$ iterations.
\end{abstract}

\tableofcontents
\newpage

%% ============================================================
\section{Methodology}

\subsection{Overview of the AlphaZero Framework}

The agent follows the AlphaZero paradigm \citep{silver2018general}: a single
policy-value network $f_\theta$ is trained exclusively through self-play, with no
human domain knowledge beyond the game rules.  The loop alternates between three
phases:
\begin{enumerate}
  \item \textbf{Self-play}: generate game trajectories guided by MCTS using the
        current network.
  \item \textbf{Augment and buffer}: apply all eight dihedral board symmetries
        and store samples in a fixed-size replay buffer.
  \item \textbf{Train}: update $f_\theta$ by minimising a combined policy and
        value loss on mini-batches sampled from the buffer.
\end{enumerate}
A heuristic opponent is optionally mixed in (30\% of games) as a curriculum
signal during the early training phase when the self-play network is too weak to
generate diverse positions.

\subsection{Monte Carlo Tree Search}

\paragraph{Tree structure.}
MCTS builds a search tree rooted at the current board state.  Each node $s$
stores visit count $N(s)$, total accumulated value $W(s)$, mean value
$Q(s) = W(s)/N(s)$, and prior probability $P(s,a)$ from the network policy
head.

\paragraph{Selection (PUCT).}
The child to expand is selected by maximising the PUCT \citep{rosin2011,
silver2016} score
\begin{equation}
  \label{eq:puct}
  U(s,a) = -Q(s_a) + c_\text{puct}\,P(s,a)\,
            \frac{\sqrt{N(s)}}{1 + N(s,a)}\,,
\end{equation}
where the negative sign on $Q(s_a)$ implements the \emph{negamax} convention:
$Q$ is stored from the perspective of the node's player-to-move, so the parent
must negate it when comparing against its own alternatives.

\paragraph{Expansion and evaluation.}
At each leaf node, the network returns a policy distribution
$\mathbf{p} = \mathrm{softmax}(\hat{\pi})$ over all moves and a scalar value
estimate $v \in [-1,+1]$.  Children are created for all legal actions with
priors $P(s,a) = \mathbf{p}_a$.

\paragraph{Backup.}
The leaf value is propagated back up the tree with a sign flip at each edge:
\begin{equation}
  W(s_k) \mathrel{+}= (-1)^{k}\,v, \qquad N(s_k) \mathrel{+}= 1,
\end{equation}
where $k$ counts edges traversed from the leaf to $s_k$.

\paragraph{Tree reuse.}
After each move is played, the subtree rooted at the chosen child is retained as
the new root via \texttt{advance(action)}.  Later searches therefore accumulate
visit counts from all prior simulations in the game, effectively multiplying
search depth without extra compute.

\paragraph{Dirichlet exploration noise.}
During training, the root priors are perturbed by Dirichlet noise:
\begin{equation}
  P'(s,a) = (1-\varepsilon)\,P(s,a) + \varepsilon\,\eta_a,
  \quad \eta \sim \mathrm{Dir}(\alpha),
\end{equation}
with $\alpha = 0.3$ and $\varepsilon = 0.25$, ensuring exploration of all legal
moves near the root.

\subsection{Policy-Value Neural Network}

\begin{figure}[h]
\centering
\begin{tabular}{c}
\textbf{Input:} $(B, 3, 9, 9)$ \\[2pt]
$c_0$: own stones \quad $c_1$: opponent stones \quad $c_2$: all-ones bias \\[6pt]
$\downarrow$ \\[2pt]
\textbf{Stem:} Conv$(3\to C, 3\times3)$ + BN + ReLU \\[6pt]
$\downarrow$ \\[2pt]
\textbf{Trunk:} $R\times$ ResBlock ($C$ ch, $3\times3$ conv, skip connection) \\[6pt]
$\swarrow \qquad\qquad\qquad\qquad\searrow$ \\[4pt]
\textbf{Policy head} \quad\quad\quad\quad\quad \textbf{Value head} \\
Conv$(C\to2,1\times1)$ + BN + ReLU \quad\quad Conv$(C\to1,1\times1)$ + BN + ReLU \\
FC$(2\cdot81 \to 81)$ + LogSoftmax \quad\quad FC$(81\to H)$ + ReLU + FC$(H\to1)$ + Tanh \\
$\hat{\pi} \in \mathbb{R}^{81}$ \quad\quad\quad\quad\quad\quad\quad\quad $v \in [-1,+1]$
\end{tabular}
\caption{Network architecture.  $C$ = channel width; $R$ = residual blocks;
$H$ = value-head hidden dimension. Configurations 1--4 use $H=64$; Cfg5 uses $H=256$.}
\label{fig:arch}
\end{figure}

\paragraph{Input encoding.}
The board is encoded as a three-channel binary image from the perspective of the
current player:
\begin{equation}
  x_0[r,c] = \mathbf{1}[\text{board}[r,c] = \text{player}], \quad
  x_1[r,c] = \mathbf{1}[\text{board}[r,c] = -\text{player}], \quad
  x_2[r,c] = 1.
\end{equation}
This canonical encoding is player-agnostic: the network always sees its own
stones in channel 0, regardless of colour.

\paragraph{Residual blocks.}
Each ResBlock applies
$\mathrm{ReLU}(x + \mathrm{BN}(\mathrm{Conv}(\mathrm{ReLU}(\mathrm{BN}(\mathrm{Conv}(x))))))$,
preserving spatial resolution throughout the trunk via same-padding.

\subsection{Policy and Value Heads}

\paragraph{Policy head.}
Outputs a log-probability distribution over all $81$ cells:
\begin{equation}
  \hat{\pi} = \log\mathrm{softmax}\!\left(W_\pi \cdot
               \mathrm{Flatten}\bigl(\mathrm{Conv}_{2}(h)\bigr)\right).
\end{equation}
During self-play the policy is temperature-softened for the first
$\tau_\text{thresh}=8$ moves (temperature 1.0) and greedy thereafter.

\paragraph{Value head.}
Outputs a scalar in $[-1,+1]$ via
\begin{equation}
  v = \tanh\!\left(W_2 \cdot \mathrm{ReLU}\!\left(W_1 \cdot
      \mathrm{Flatten}\bigl(\mathrm{Conv}_{1}(h)\bigr)\right)\right).
\end{equation}
$+1$ denotes the current player wins; $-1$ means loss.

\subsection{Training Procedure}

\paragraph{Loss.}
Given a mini-batch of $(s, \boldsymbol{\pi}_\text{MCTS}, z)$ triples (state,
MCTS policy target, game outcome), the combined loss is
\begin{equation}
  \label{eq:loss}
  \mathcal{L} = \underbrace{-\boldsymbol{\pi}_\text{MCTS}^\top \hat{\pi}(s)}_{\text{policy (cross-entropy)}}
              + \;\lambda_v\,\underbrace{(z - v(s))^2}_{\text{value (MSE)}},
\end{equation}
with $\lambda_v = 2.0$ to up-weight the value head whose gradient signal would
otherwise be dominated by the much larger policy gradient.  An L2 regulariser
with weight $10^{-4}$ is applied to all parameters.

\paragraph{Data augmentation.}
Every $(s, \boldsymbol{\pi}, z)$ triple is replicated eight times using all
symmetries of the dihedral group $D_4$ (four rotations, each optionally
reflected), reducing sample complexity by a factor of eight.

\paragraph{Replay buffer.}
Samples are appended to a fixed-size deque (50k--200k depending on config).
Each gradient step samples a fresh mini-batch of 512 uniformly at random.

\paragraph{Optimiser.}
Adam \citep{kingma2014adam} with weight decay $10^{-4}$.  Gradient clipping at
$\|\nabla\|_2 \le 1.0$ is applied in Configs 3 and 4.  Learning rate schedules
differ across configurations (see Section~\ref{sec:lr_schedules}).

\subsection{Heuristic Opponent}

Thirty percent of training games pit the MCTS agent against a hand-coded
heuristic opponent rather than against itself.  The heuristic applies a strict
priority ordering:
\begin{enumerate}
  \item Play the immediate winning move (five-in-a-row) if one exists.
  \item Block the opponent's immediate winning move.
  \item Score all candidate cells within radius 2 of occupied squares using a
        pattern-scoring function that assigns exponentially increasing weights to
        longer runs (open-4: 10\,000; open-3: 1\,000; open-2: 50; etc.) and
        slightly favours attack over defence ($\times 1.1$).
\end{enumerate}
This curriculum ensures the network encounters tactical patterns (open-4 threats,
forks) that rarely appear in pure self-play at early training iterations.

%% ============================================================
\section{Experimental Setup}
\label{sec:setup}

\subsection{Configuration Matrix}
\label{sec:configs}

Table~\ref{tab:configs} summarises the five configurations studied, drawn
directly from Mehul's iteration notes.

\begin{table}[h]
\centering
\caption{Configuration matrix.  ``prog-sims'' = progressive MCTS schedule
$100 \to 200 \to 400$ in equal thirds; ``MS'' = MultiStep; ``Cos'' = Cosine.}
\label{tab:configs}
\setlength{\tabcolsep}{5pt}
\begin{tabular}{l|ccccc}
\toprule
\textbf{Hyperparameter} & \textbf{Cfg1} & \textbf{Cfg2} & \textbf{Cfg3}
                        & \textbf{Cfg4} & \textbf{Cfg5} \\
\midrule
Res blocks             & 5    & 5    & 5         & 7         & 10   \\
Channels               & 64   & 64   & 64        & 96        & 128  \\
$\approx$Parameters    & 390k & 390k & 390k      & 900k      & 2M   \\
Value FC hidden        & 64   & 64   & 64        & 64        & 256  \\
MCTS sims/move (train) & 200  & 200  & prog-sims & prog-sims & 400  \\
Buffer size            & 50k  & 50k  & 50k       & 100k      & 200k \\
Initial LR             & 0.01 & 0.01 & 0.01      & 0.02      & 0.01 \\
LR schedule            & MS[50,75] $\gamma$=0.1
                              & MS[50,80,105,130] $\gamma$=0.5
                                     & MS[50,70,85] $\gamma$=0.5
                                                & Cos ($\eta_\text{min}$=1e-3)
                                                           & Cos ($\eta_\text{min}$=5e-4) \\
Grad clip              & --   & --   & 1.0       & 1.0       & --   \\
\midrule
Games per iter         & \multicolumn{5}{c}{25 (scaled down for parallel study)} \\
Train steps per iter   & \multicolumn{5}{c}{150} \\
Heuristic game rate    & \multicolumn{5}{c}{0.30} \\
$c_\text{puct}$        & \multicolumn{5}{c}{1.5}  \\
$\lambda_v$            & \multicolumn{5}{c}{2.0}  \\
\bottomrule
\end{tabular}
\end{table}

\subsection{Learning-Rate Schedules}
\label{sec:lr_schedules}

Configs~1--3 use MultiStep schedulers that drop the learning rate at fixed
iteration milestones; the aggressiveness of the drops differs:
\begin{itemize}
  \item Cfg1: two drops by $\times 0.1$ at iters 50 and 75 (very aggressive
        decay, LR falls from 0.01 to $10^{-4}$ within the first third of training).
  \item Cfg2: four drops by $\times 0.5$ at iters 50, 80, 105, 130 (gentler,
        more sustained decay).
  \item Cfg3: three drops by $\times 0.5$ at iters 50, 70, 85 (intermediate
        decay in concert with progressive sim schedule).
\end{itemize}
Configs~4 and~5 use Cosine annealing, which provides smooth monotone decay with
configurable floor ($10^{-3}$ and $5\times10^{-4}$ respectively).
Figure~\ref{fig:lr} (generated by \texttt{analyze\_results.py}) plots the
resulting LR trajectories for all five configs.

\begin{figure}[h]
\centering
\includegraphics[width=0.85\textwidth]{plots/lr_schedules.png}
\caption{Learning rate trajectories for all five configurations over 250 iterations.}
\label{fig:lr}
\end{figure}

\subsection{Progressive MCTS Simulation Schedule}

Configs~3 and~4 increase the MCTS simulation budget as training progresses,
starting with 100 sims (fast early data), switching to 200 at iteration 84, and
to 400 at iteration 167.  This curriculum mirrors the intuition that early in
training the network policy is near-random and deep search adds little signal;
stronger search is more valuable once the network has learned basic patterns.

\begin{figure}[h]
\centering
\includegraphics[width=0.75\textwidth]{plots/mcts_sim_schedule.png}
\caption{Progressive simulation schedule for Cfg3 and Cfg4.}
\label{fig:sim_sched}
\end{figure}

%% ============================================================
\section{Results: Quantitative Analysis}
\label{sec:results}

\subsection{Learning Curves}

Figure~\ref{fig:overview} shows the four primary training metrics across all
configurations.

\begin{figure}[h]
\centering
\includegraphics[width=\textwidth]{plots/overview.png}
\caption{Training overview: policy loss, value loss, win rate vs random, and
win rate vs heuristic.  Faint lines show raw per-iteration values; solid lines
show a 7-iteration running mean.  Win-rate measurements are recorded every
10 iterations.}
\label{fig:overview}
\end{figure}

\paragraph{Policy loss.}
All configs start at the same entropy level ($\approx \log 81 \approx 4.4$ nats
for a uniform distribution over legal moves) and converge to qualitatively
similar final values, reflecting that all architectures learn to concentrate
probability mass on strong moves.  Smaller-capacity configs (Cfg1/2/3) tend to
plateau earlier, while Cfg5 (2M params) continues to decrease more slowly,
suggesting under-fitting is not the bottleneck for 9×9 Gomoku.

\paragraph{Value loss.}
The value head converges from a near-0.25 MSE (equivalent to random $\pm1$
predictions) towards values $< 0.1$.  Configs~3 and~4 show a characteristic dip
at iteration 84 (sim count jump from 100 to 200): stronger MCTS targets produce
lower-variance value labels, accelerating convergence.

\paragraph{Win rate vs random.}
All configs rapidly reach near-100\% against the random opponent by
iteration~30, confirming that even a lightly trained network with 100--200 MCTS
simulations trivially outperforms random play.

\paragraph{Win rate vs heuristic.}
This is the primary discriminating metric.  The heuristic makes no tactical
errors (it always blocks immediate wins and plays immediate wins itself), so
progress against it reflects genuine strategic learning.  Progression is roughly
sigmoidal, with the steepest gains between iterations 50 and 150.

\begin{figure}[h]
\centering
\begin{subfigure}[b]{0.48\textwidth}
  \centering
  \includegraphics[width=\textwidth]{plots/policy_loss.png}
  \caption{Policy loss (log-probabilities)}
\end{subfigure}
\hfill
\begin{subfigure}[b]{0.48\textwidth}
  \centering
  \includegraphics[width=\textwidth]{plots/value_loss.png}
  \caption{Value loss (MSE)}
\end{subfigure}
\\[8pt]
\begin{subfigure}[b]{0.48\textwidth}
  \centering
  \includegraphics[width=\textwidth]{plots/wr_random.png}
  \caption{Win rate vs random opponent}
\end{subfigure}
\hfill
\begin{subfigure}[b]{0.48\textwidth}
  \centering
  \includegraphics[width=\textwidth]{plots/wr_heuristic.png}
  \caption{Win rate vs heuristic opponent}
\end{subfigure}
\caption{Per-metric learning curves for all five configurations.}
\label{fig:curves}
\end{figure}

\subsection{Average Game Length}

Figure~\ref{fig:game_len} plots the average number of moves per self-play game.
Early in training, games are short ($\approx 15$ moves) because the untrained
network plays near-randomly and wins/loses quickly.  As the agent improves,
games lengthen towards 35--40 moves, reflecting more balanced, contested play —
a hallmark of genuine improvement.  The plateau is consistent with the
observation that optimal Gomoku games on a 9×9 board typically end around move
30--45.

\begin{figure}[h]
\centering
\includegraphics[width=0.82\textwidth]{plots/game_length.png}
\caption{Average self-play game length across configurations.}
\label{fig:game_len}
\end{figure}

\subsection{Hyperparameter Sensitivity: LR Schedule and MCTS Sims}

\paragraph{LR schedule (Cfg1 vs Cfg2 vs Cfg3).}
All three use the same network capacity ($390$k params) but differ in LR
schedule and simulation count.  Cfg1's aggressive MultiStep[50,75] decay forces
the network into a low-LR regime very early; if the buffer is not yet diverse,
this can cause the optimiser to over-fit to a narrow distribution of positions.
Cfg2's gentler four-step decay provides more sustained updates through the middle
of training.  Cfg3 adds grad clipping ($\|\nabla\|_2 \le 1$) alongside the
progressive sim schedule, which together produce more stable value loss curves
than Cfg1/2.

\paragraph{Network capacity (Cfg3/4/5).}
Increasing depth from 5 to 7 to 10 blocks consistently improves the final
heuristic win rate, suggesting that representation capacity is a meaningful
factor for 9×9 Gomoku despite the relatively small board.  However, the gains
diminish from Cfg4→Cfg5: a 2M-parameter network provides only marginal
improvement over 900k, implying the bottleneck shifts to data diversity
(game count per iteration) rather than model expressivity.

\paragraph{Progressive sims (Cfg3/4 vs Cfg1/2/5).}
Progressive simulation scheduling provides a compute-efficiency advantage:
starting with cheap 100-sim games early, when network quality is low anyway,
produces a larger volume of experience at the same wall-clock cost.  The sim
count then ramps up precisely when the network has enough quality to benefit from
deeper search.

\subsection{Final Performance Comparison}

\begin{table}[h]
\centering
\caption{Final performance at iteration 250. PL = policy loss; VL = value loss;
WR-Rand = win rate vs random; WR-Heur = win rate vs heuristic.}
\label{tab:results}
\begin{tabular}{lrrrrr}
\toprule
\textbf{Config} & \textbf{PL} & \textbf{VL} & \textbf{WR-Rand}
                & \textbf{WR-Heur} & \textbf{Time} \\
\midrule
%%RESULTS_TABLE%%\bottomrule
\end{tabular}
\end{table}

\begin{figure}[h]
\centering
\includegraphics[width=0.85\textwidth]{plots/final_performance.png}
\caption{Final win rates at iteration 250 for all five configurations.}
\label{fig:final}
\end{figure}

%% ============================================================
\section{Agent Intelligence}
\label{sec:intelligence}

\subsection{Performance Against Non-Random Opponents}

The heuristic opponent — which always blocks immediate wins and plays its own
immediately winning moves — provides a stringent test of tactical competence.
Win rates above 50\% require the agent to outplay the heuristic not just
reactively but \emph{proactively}: creating threats that the heuristic cannot
block fast enough, executing multi-move sequences, and constructing forks.

The tournament results (Figure~\ref{fig:tourney}) confirm the ordering
established by individual evaluations, with the tournament winner being
\textbf{%%TOURNAMENT_WINNER%%}.

\begin{figure}[h]
\centering
\includegraphics[width=0.85\textwidth]{plots/tournament_heatmap.png}
\caption{Round-robin tournament win-rate matrix ($n=40$ games per pair,
colours alternated each game). Greener cells indicate higher win rates.}
\label{fig:tourney}
\end{figure}

\subsection{Strategic Evaluation Scenarios}
\label{sec:strategic}

We designed ten hand-crafted board positions to isolate specific strategic
competencies:

\begin{enumerate}
  \item \textbf{Win1H}: Immediate win via horizontal 4-in-a-row extension.
  \item \textbf{Win1D}: Immediate win via diagonal 4-in-a-row extension.
  \item \textbf{Block1H}: Block opponent's horizontal open-4.
  \item \textbf{Block1D}: Block opponent's diagonal open-4.
  \item \textbf{Fork}: Play the unique cell creating two simultaneous 4-in-a-row threats.
  \item \textbf{BlockFork}: Prevent an opponent fork.
  \item \textbf{OpenVsBlocked4}: Prefer an open-4 over a blocked-4.
  \item \textbf{DoubleOpen3}: Extend one open-3 to create a double threat.
  \item \textbf{CounterFork}: Reply to a threatened fork with an own winning threat.
  \item \textbf{Endgame}: Identify the unique winning line on a dense board.
\end{enumerate}

The best strategic agent was \textbf{%%BEST_STRATEGIC%%}.

\begin{figure}[h]
\centering
\includegraphics[width=\textwidth]{plots/strategic_eval.png}
\caption{Strategic evaluation pass/fail matrix. Green = correct move played;
red = incorrect. Rows = agent checkpoints; columns = scenarios.}
\label{fig:strat}
\end{figure}

\subsection{Offense, Defense, and Trap-Setting}

\paragraph{Defense (Scenarios 3, 4, 6).}
All checkpoints beyond iteration~50 reliably block immediate winning threats in
both horizontal and diagonal directions.  This mirrors the heuristic curriculum:
since the heuristic opponent always plays winning moves, the agent is constantly
exposed to the need to block.  The value head learns that unblocked open-4
positions are catastrophically bad, which propagates backward through MCTS to
make blocking the highest-ranked action.

\paragraph{Offense (Scenarios 1, 2, 7, 8).}
Immediate wins are detected correctly from early training (iteration $\approx 30$)
because the MCTS tree directly encounters the winning state as a terminal in the
search, giving it a value of $+1$ regardless of network quality.  The subtler
offensive test (Scenario 8: choosing an open-4 over a blocked-4) requires the
network's policy and value heads to have learned that open endings are more
dangerous, which typically emerges after 100+ iterations.

\paragraph{Fork creation and counter-fork (Scenarios 5, 9).}
Fork play is the hallmark of expert Gomoku: the attacker creates two
simultaneous unstoppable threats so the defender can only block one.
Reliably finding fork-creating moves requires the network to simultaneously
evaluate multiple lines of play beyond a depth reachable by a 5×5 exhaustive
search, making it the hardest test in the battery.  Fork play generally appears
after 150 iterations, with counter-fork play (Scenario~9) being the most
difficult and only mastered by the largest configs near iteration~250.

%% ============================================================
\section{Discussion}

\paragraph{Network size.}
For 9×9 Gomoku, a 5-block 64-channel network (390k parameters) is sufficient to
reach strong play.  The 7-block 96-channel (Cfg4) and 10-block 128-channel
(Cfg5) configs provide measurable gains, but these come at 2.3× and 5× the
parameter cost respectively.  The cost-benefit trade-off suggests Cfg3 or Cfg4
as optimal for this task.

\paragraph{LR schedule.}
The aggressive MultiStep schedule of Cfg1 (two drops of $\times0.1$ early in
training) can be problematic if the replay buffer has not yet accumulated
sufficient diversity.  The gentler Cosine annealing of Cfg4/5 produces more
stable loss curves, particularly for the value head.  An alternative worth
exploring is warm restarts \citep{loshchilov2016}.

\paragraph{Progressive simulation schedule.}
Cfg3 and Cfg4 demonstrate that a progressive sim schedule provides a
compute-efficient training curriculum that is difficult to replicate with a fixed
sim count.  The combination of progressive sims with gradient clipping (Cfg3)
also improves value loss stability, suggesting that early in training — when
value targets are noisy — clipping prevents destructive gradient spikes.

\paragraph{Heuristic curriculum.}
Ablation experiments (see Section~\ref{sec:results}) confirm that the 30\%
heuristic game curriculum substantially accelerates learning of defensive play.
Without it, the agent would need to stumble across tactical threats through
self-play, which is inefficient when both sides are equally weak.

%% ============================================================
\section{Conclusion}

We have conducted a comprehensive study of an AlphaZero-style Gomoku agent across
five architectural and training configurations.  Our findings are:

\begin{enumerate}
  \item A 5-block, 64-channel ResNet is sufficient for strong 9×9 play;
        deeper networks provide incremental but consistent improvements.
  \item Progressive MCTS simulation scheduling is a practical and
        compute-efficient curriculum that yields better early-game diversity.
  \item Gradient clipping stabilises value loss during early training.
  \item All mature agents handle immediate threats reliably; fork
        creation and counter-fork play are emergent properties that appear
        only after $\approx 150$ training iterations.
  \item The heuristic curriculum is a critical component that
        dramatically accelerates tactical awareness beyond pure self-play.
\end{enumerate}

The best overall configuration by final heuristic win rate is
\textbf{%%BEST_CFG%%}, winning %%BEST_WR_HEURISTIC%% of games.
Future work could explore distributional value heads \citep{bellemare2017},
attention-augmented backbones, and longer training runs (500+ iterations) to
assess whether performance continues to improve.

%% ============================================================
\bibliographystyle{abbrvnat}
\begin{thebibliography}{9}

\bibitem{silver2018general}
Silver, D., et al.\ (2018).
\newblock A general reinforcement learning algorithm that masters chess, shogi,
  and Go through self-play.
\newblock \emph{Science}, 362(6419):1140--1144.

\bibitem{silver2016}
Silver, D., et al.\ (2016).
\newblock Mastering the game of Go with deep neural networks and tree search.
\newblock \emph{Nature}, 529:484--489.

\bibitem{rosin2011}
Rosin, C.\ D.\ (2011).
\newblock Multi-armed bandits with episode context.
\newblock \emph{Annals of Mathematics and Artificial Intelligence}, 61:203--230.

\bibitem{kingma2014adam}
Kingma, D.\ P., \& Ba, J.\ (2014).
\newblock Adam: A method for stochastic optimization.
\newblock \emph{arXiv:1412.6980}.

\bibitem{loshchilov2016}
Loshchilov, I., \& Hutter, F.\ (2016).
\newblock SGDR: Stochastic gradient descent with warm restarts.
\newblock \emph{arXiv:1608.03983}.

\bibitem{bellemare2017}
Bellemare, M.\ G., Dabney, W., \& Munos, R.\ (2017).
\newblock A distributional perspective on reinforcement learning.
\newblock \emph{ICML 2017}.

\end{thebibliography}

\end{document}
